\begin{table*}[!t]
\centering
\caption{Module~8 constants (b): CRAC, the relay output edge, pools and diffusion constants.}
\label{tab:m8constb}
\footnotesize
\renewcommand{\arraystretch}{1.15}
\begin{tabularx}{\textwidth}{@{}l r l c >{\raggedright\arraybackslash}X@{}}
\toprule
Symbol & Value & Unit & Basis & Meaning, justification, source\\
\midrule
$k_\mathrm{Cr}$ & \num{10} & $\mu\mathrm{M}^{-1}\mathrm{s}^{-1}$ & C & Phosphorylation of membrane-bound CRAC by either PKB. \mkt{M8-tor-pkbs}{Both PKBs are TORC2 targets}~\cite{kamimura2008} and \mkt{M8-aca-prolong}{RasC--TORC2--PKBR1 signalling prolongs ACA activation}~\cite{cai2010}; PKB$\to$CRAC is a model assumption. \\
$k_{\mathrm{Cr},P}$ & \num{1} & $\mu\mathrm{M}^{-1}\mathrm{s}^{-1}$ & E & PIP3-facilitated recruitment of CRAC, which is \mkt{M8-crac-pip3}{a PH-domain protein}~\cite{insall1994,parent1998,comer2005}. \\
$k_{\mathrm{Cr},b}$ & \num{0.2} & $\mathrm{s^{-1}}$ & C & PIP3-independent membrane association; keeps the relay functional when PIP3 is absent (\mkt{M8-pi3k-ind}{PI3K-independent PKB signalling}~\cite{kamimura2008}). \\
$k_{\mathrm{Cr},\mathrm{off}}$ & \num{0.1} & $\mathrm{s^{-1}}$ & E & Release of unphosphorylated CRAC from the membrane. \\
$k_{\mathrm{Cr},\mathrm{dp}}$ & \num{1.5} & $\mathrm{s^{-1}}$ & D & $=k_\mathrm{Cr}\,a\,(1-f)/f$ with $a=567$ PKBA$^*$/voxel (the measured peak) and target CRAC$^*$ fraction $f=0.44$; fixes the activated fraction of the CRAC switch. \\
$k_X$ & \num{200} & $\mathrm{molec\,s^{-1}\,cell^{-1}}$ & C & Injection rate of X at full CRAC$^*$ (per-CRAC$^*$ rate $k_X/N_\mathrm{CRAC}=200/2002=0.1\,\mathrm{s^{-1}}$, with $N_\mathrm{CRAC}=2002$ per cell); must exceed the $\approx\!24$-molecule threshold of Module~0. \\
$K_P,\ n_P$ & \num{5000},\ 8 & $\mathrm{molec/voxel},\ -$ & C & Half-open level and steepness of the PIP3 gate (Sec.~\ref{sec:hill}). \\
$k_\mathrm{PKA}$ & \num{0.938} & $\mu\mathrm{M}^{-1}\mathrm{s}^{-1}$ & E & PKA$^*$ feedback on PKBR1$^*$; \mkt{M8-pka-fb}{PKA exerts negative feedback on Ras, Rap1 and TORC2 signalling}~\cite{scavello2017}. The magnitude is a design choice: the value $0.6\,\mathrm{min^{-1}}$ is \mkt{M8-laub-k}{the coefficient $k_5$ of}~\cite{laub1998} (CAR1$\to$ERK2, not a PKA coefficient; the PKA$\to$ERK2 coefficient is $k_6=0.8\,\mathrm{min^{-1}}$, used in Module~10), $30\times$ boosted and normalised by the PKA pool. \\
$[\mathrm{RasC}]_\mathrm{tot}$ & \num{385} & $\mathrm{molec/voxel}$ & E & Pool size; not constrained by measurements. \\
$[\mathrm{CGAP}]_\mathrm{tot},\ [\mathrm{TORC2}]_\mathrm{tot},\ [\mathrm{CRAC}]_\mathrm{tot}$ & \num{154} & $\mathrm{molec/voxel}$ & E & Pool sizes (2\,002 per cell); not constrained by measurements. \\
$[\mathrm{PKBR1}]_\mathrm{tot}$ & \num{77} & $\mathrm{molec/voxel}$ & E & PKBR1 pool; PkbA (Module~4) is a separate, larger pool. \\
$D_{\mathrm{RasC}},D_{\mathrm{TORC2}},D_{\mathrm{PKBR1}},D_{\mathrm{CRAC}_m}$ & \num{0.1} & $\mu\mathrm{m^2/s}$ & E & Membrane-bound species. \\
$D_{\mathrm{GEF_A}},D_{\mathrm{CGAP}},D_{\mathrm{CRAC}_c}$ & \num{10} & $\mu\mathrm{m^2/s}$ & E & Cytosolic value ($D_{\mathrm{GEF_A^*}}=2$). \\
\bottomrule
\end{tabularx}
\end{table*}
